\section{Morphology}\label{sec:morphology}

All inflection is suffixal, agglutinative and stress-neutral. Every ending is itself the product of a family vote recorded in the Manual, so the grammar is derived by the same procedure as the lexicon.

\subsection{Nouns}\label{sec:nouns}
\paragraph{No gender.} Finnic has none; Baltic, Slavic and Germanic disagree with one another about which nouns are what. Gender assignment cannot be made deterministic across the nine and is abolished. This is also the honest outcome of the vote: three gendered families with three incompatible systems cannot agree on a form.

\paragraph{No articles.} Finnic, Baltic and Slavic have none: three votes to one.

\paragraph{Plural \emph{-i}.} Baltic (lv \emph{-i}, lt \emph{-ai/-iai}) and Slavic (\emph{-i/-y}) agree on the vowel; Finnic \emph{-t/-d} and Germanic \emph{-er/-e/-en} each cast one vote: 2--1--1. \emph{kala} $\rightarrow$ \emph{kalai}, \emph{hülje} $\rightarrow$ \emph{hüljei}, \emph{sort} $\rightarrow$ \emph{sorti}.

\paragraph{Seven cases.} Finnic's rich case system and the Baltic--Slavic middle ground produce seven cases (Table~\ref{tab:cases}); Germanic contributes little because its cases are mostly gone. Plural and case stack, stem + \emph{-i} + case: \emph{kalju-i-l} ``on the rocks'', \emph{Baltmeri-se} ``in the Baltic''. The accusative plural is identical to the nominative plural, as for inanimates in Slavic and generally in Finnic and Germanic. Compounds are head-final and written solid, as in every source language: \emph{peleekhülje} ``grey seal'' (the species) against \emph{peleek hülje} ``a seal that happens to be grey''.

\begin{table}[t]
\centering\small
\caption{Noun cases and the votes that produced them.}
\label{tab:cases}
\setlength{\tabcolsep}{4pt}
\begin{tabular}{@{}llp{1.8cm}p{2.9cm}@{}}
\toprule
Case & Ending & Meaning & Source of the vote \\
\midrule
Nominative & -Ø & subject & unanimous \\
Genitive & -n & of, 's & Finnic -n, Germanic -s, Slavic -a tie $\rightarrow$ Estonian \\
Accusative & -a & direct object & all abstain $\rightarrow$ Fallback $\rightarrow$ Lithuanian -ą \\
Inessive & -se & in, inside & only Finnic has a suffix (-ssa/-s) \\
Adessive & -l & on, at; also \emph{when} & Finnic -lla/-l \\
Elative & -ste & from, out of & Finnic -sta/-st \\
Allative & -le & to, onto & Finnic -lle/-le \\
\bottomrule
\end{tabular}
\end{table}

\subsection{Adjectives, comparison, adverbs}
Adjectives are bare stems and precede the noun (all nine put adjectives first). They agree in number only---Germanic-style partial agreement, since the full-agreement families cannot agree on \emph{how}: \emph{duž kala} ``big fish'', \emph{duži kalai} ``big fish (pl.)''. Comparative \emph{-m} (Germanic \emph{-er/-are}, Slavic \emph{-ejš/-sz}, Finnic \emph{-mpi/-m}, three-way tie $\rightarrow$ Estonian): \emph{dužim} ``bigger''. Superlative \emph{-st} (Germanic \emph{-st}, Slavic \emph{naj-} tie $\rightarrow$ Swedish): \emph{pospolitist} ``most common''. Adverb \emph{-t} (Finnic and Germanic agree on the dental, 2--1--1): \emph{dužit} ``greatly''. Derived adjective \emph{-ne} (Finnic \emph{-(i)nen/-ne}, Slavic \emph{-ny}): \emph{solne} ``salty'', \emph{zimne} ``cold''.

\subsection{Pronouns and numerals}\label{sec:pronouns}
Table~\ref{tab:pronouns} lists the personal pronouns and the numerals one to five, each with its vote. Genitives are \emph{jan, tun, hanen, men, juusen, den}. The third person is genderless, like Finnish \emph{hän}. The numerals alone show the Rota Rule at work: five words, five different tie-break winners (Finnish, plurality, plurality, Baltic--Slavic, Danish). Six to ten are derived by the same procedure and have no manual decree: \emph{seši, septini, ota, dzevents, desmit}.

\begin{table}[t]
\centering\small
\caption{Pronouns and lower numerals with their votes.}
\label{tab:pronouns}
\begin{tabular}{@{}llp{4.6cm}@{}}
\toprule
& Baltmeri & Vote \\
\midrule
I & ja & pl/ru \emph{ja}, sv \emph{jag}, da \emph{jeg}: J in two families \\
you & tu & T-onset in three families; \emph{u} beats \emph{y} \\
he/she/it & han & fi \emph{hän} + sv/da \emph{han}: H-N in two families \\
we & me & M-onset in three families; Baltic final -s dropped \\
you (pl.) & juus & Finnic \emph{te}, Baltic \emph{jūs}, Slavic \emph{vy} tie $\rightarrow$ Latvian \\
they & de & Slavic \emph{oni}, Germanic \emph{de} tie $\rightarrow$ Swedish \\
\midrule
1 & üks & four-way tie $\rightarrow$ Finnish \emph{yksi} \\
2 & dva & D-V in Baltic, Slavic, Germanic \\
3 & tri & T-R everywhere except Finnic \\
4 & četri & Baltic + Slavic, one differing position \\
5 & fem & Slavic vs Germanic tie $\rightarrow$ Danish \\
\bottomrule
\end{tabular}
\end{table}

\subsection{Verbs}
The infinitive is \emph{-te}: Baltic \emph{-t/-ti} and Slavic \emph{-ć/-t'} agree on the dental, plus the supporting \emph{-e}. Finnic, Baltic and Slavic all inflect for person; only Swedish and Danish do not, so 3--1, Baltmeri inflects. Each slot was voted separately (Table~\ref{tab:verbs}). The third person is unmarked for number, Scandinavian-style; the plural subject carries the plural. Past \emph{-l-} (Germanic dental vs Slavic \emph{-l} tie $\rightarrow$ Polish) and future \emph{-s-} (Baltic is the only family with a synthetic future, and the Manual prefers a suffix to a helper verb wherever one exists) are inserted between stem and person ending. Negation is \emph{ne} before the verb (Baltic \emph{ne-}, Slavic \emph{nie/ne}); after the copula, \emph{ne} precedes the predicate (\emph{Baltmeri er ne glembok meri} ``the Baltic is not a deep sea''). The answer-word ``no'' is \emph{nei} (Germanic \emph{nej} + Finnic \emph{ei}). The copula is suppletive, as in all nine languages (Table~\ref{tab:copula}): present \emph{es-} is a Hanse word (lv \emph{esmu}, lt \emph{esu}, pl \emph{jestem}); the third person \emph{er} is a separate vote (lv \emph{ir}, lt \emph{yra}, sv \emph{är}, da \emph{er}).

\begin{table}[t]
\centering\small
\caption{Person endings and the full regular paradigm of \emph{vidte} ``to see''. Forms in italics show Repair (\S\ref{sec:repair}).}
\label{tab:verbs}
\begin{tabular}{@{}llp{2.6cm}@{}}
\toprule
& Ending & Vote \\
\midrule
1sg & -u & lv, lt, ru \emph{-u}: two families \\
2sg & -t & Finnic \emph{-t/-d}, Baltic \emph{-i}, Slavic \emph{-š} tie $\rightarrow$ Finnish \\
3sg & -Ø & lv, pl, sv, da all zero \\
1pl & -me & M in all four families \\
2pl & -te & T in all four families \\
3pl & -Ø & Finnic \emph{-vat} vs Germanic Ø tie $\rightarrow$ Danish \\
\bottomrule
\end{tabular}

\medskip
\begin{tabular}{@{}llll@{}}
\toprule
& Present & Past (-l-) & Future (-s-) \\
\midrule
ja & vidu & vidlu & vidsu \\
tu & \emph{vidit} & \emph{vidlit} & \emph{vidsit} \\
han / de & vid & \emph{vidil} & vids \\
me & vidme & \emph{vidilme} & \emph{vidisme} \\
juus & vidte & \emph{vidilte} & \emph{vidiste} \\
\bottomrule
\end{tabular}
\end{table}

\begin{table}[t]
\centering\small
\caption{The copula \emph{bute} ``to be''.}
\label{tab:copula}
\begin{tabular}{@{}llll@{}}
\toprule
& Present & Past & Future \\
\midrule
ja & esu & bulu & busu \\
tu & est & bulit & bust \\
han / de & er & bul & bus \\
me & esme & bulme & busme \\
juus & este & bulte & buste \\
\bottomrule
\end{tabular}
\end{table}

\subsection{Derivation}
Four productive suffixes were voted: \emph{-us} abstract noun (Finnic \emph{-us}, Baltic \emph{-um(a)s}): \emph{dužus} ``size''; \emph{-ja} agent (Finnic \emph{-ja}, Baltic \emph{-ējs/-ėjas}): \emph{kalaja} ``fisher'', \emph{žeglja} ``sailor''; \emph{-ka} diminutive (Slavic \emph{-ek/-ka}, Estonian \emph{-ke}): \emph{kalaka} ``little fish''; \emph{-ne} adjective from noun: \emph{solne} ``salty''. Derived stems inflect like their base: \emph{dužusi} ``sizes'', \emph{kalajai} ``fishers''.
